\documentclass[12pt,a4paper]{ctexbook}
\usepackage{geometry}
\geometry{margin=2.2cm}
\usepackage{setspace}
\setstretch{1.35}
\usepackage{hyperref}
\title{全宋词选·ci.song.5000}
\author{古典文献整理}
\date{}
\begin{document}
\maketitle
\tableofcontents
\newpage
\section{渔父词・渔父}
\textbf{惠洪}\par\vspace{0.5em}
来往独龙冈畔路。\par
杖头落索闲家具。\par
后事前观如目睹。\par
非谶语。\par
须知一念无今古。\par
长笑老萧多病苦。\par
笑中与药皆狼虎。\par
蜡炬一枝非嘱付。\par
聊戏汝。\par
热来脱却娘生裤。\par

\section{渔父词・渔父}
\textbf{惠洪}\par\vspace{0.5em}
画饼充饥人笑汝。\par
一庵归扫南阳坞。\par
击竹作声方省悟。\par
徐回顾。\par
本来面目无藏处。\par
却望沩山敷坐具。\par
老师头角浑呈露。\par
珍重此恩逾父母。\par
须荐取。\par
堂堂密密声前句。\par

\section{渔父词・渔父}
\textbf{惠洪}\par\vspace{0.5em}
野鹤精神云格调。\par
逼人气韵霜天晓。\par
松下残经看未了。\par
当斜照。\par
苍烟风撼流泉绕。\par
闺阁珍奇徒照耀。\par
光无渗漏方灵妙。\par
活计现成谁管绍。\par
孤峰表。\par
一声月下闻清啸。\par

\section{渔父词・渔父}
\textbf{惠洪}\par\vspace{0.5em}
讲虎天华随玉麈。\par
波心月在那能取。\par
旁舍老僧偷指注。\par
回头觑。\par
虚空特地能言语。\par
归对学徒重自诉。\par
从前见解都欺汝。\par
隔岸有山横暮雨。\par
翻然去。\par
千岩万壑无寻处。\par

\section{渔父词・渔父}
\textbf{惠洪}\par\vspace{0.5em}
急雨颠风花信早。\par
枝枝叶叶春俱到。\par
何待小桃方悟道。\par
休迷倒。\par
出门无限青青草。\par
根不覆藏尘亦扫。\par
见精明树唯心造。\par
试借疑情看白皂。\par
回头讨。\par
灵云笑杀玄沙老。\par

\section{渔父词・渔父}
\textbf{惠洪}\par\vspace{0.5em}
万叠空青春杳杳。\par
一蓑烟雨吴江晓。\par
醉眼忽醒惊白鸟。\par
拍手笑。\par
清波不犯鱼吞钓。\par
津渡有僧求法要。\par
一桡为汝除玄妙。\par
已去回头知不峭。\par
犹迷照。\par
渔舟性懆都翻了。\par

\section{凤栖梧・蝶恋花}
\textbf{惠洪}\par\vspace{0.5em}
碧瓦笼晴烟雾绕。\par
水殿西偏，小立闻啼鸟。\par
风度女墙吹语笑。\par
南枝破腊应开了。\par
道骨不凡江瘴晓。\par
春色通灵，医得花重少。\par
爆暖酿寒空杳杳。\par
江城画角催残照。\par

\section{千秋岁}
\textbf{惠洪}\par\vspace{0.5em}
半身屏外。\par
睡觉唇红退。\par
春思乱，芳心碎。\par
空馀簪髻玉，不见流苏带。\par
试与问，今人秀整谁宜对。\par
湘浦曾同会。\par
手搴轻罗盖。\par
疑是梦，今犹在。\par
十分春易尽，一点情难改。\par
多少事，却随恨远连云海。\par

\section{青玉案}
\textbf{惠洪}\par\vspace{0.5em}
绿槐烟柳长亭路。\par
恨取次、分离去。\par
日永如年愁难度。\par
高城回首，暮云遮尽，目断人何处。\par
解鞍旅舍天将暮。\par
暗忆丁宁千万句。\par
一寸柔肠情几许。\par
薄衾孤枕，梦回人静，彻晓潇潇雨。\par

\section{西江月}
\textbf{惠洪}\par\vspace{0.5em}
十指嫩抽春笋，纤纤玉软红柔。\par
人前欲展强娇羞。\par
微露云衣霓袖。\par
最好洞天春晚，黄庭卷罢清幽。\par
凡心无计奈闲愁。\par
试拈花枝频嗅。\par

\section{西江月}
\textbf{惠洪}\par\vspace{0.5em}
大厦吞风吐月，小舟坐水眠空。\par
雾窗春晓翠如葱。\par
睡起云涛正涌。\par
往事回头笑处，此生弹指声中。\par
玉笺佳句敏惊鸿。\par
闻道衡阳价重。\par

\section{鹧鸪天}
\textbf{惠洪}\par\vspace{0.5em}
蜜烛花光清夜阑。\par
粉衣香翅绕团团。\par
人犹认假为真实，蛾岂将灯作火看。\par
方叹息，为遮拦。\par
也知爱处实难拚。\par
忽然性命随烟焰，始觉从前被眼瞒。\par

\section{清商怨}
\textbf{惠洪}\par\vspace{0.5em}
一段文章种性。\par
更谪仙风韵。\par
画戟丛中，清香凝宴寝。\par
落日清寒勒花信。\par
愁似海、洗光词锦。\par
后夜归舟，云涛喧醉枕。\par

\section{青玉案}
\textbf{惠洪}\par\vspace{0.5em}
凝祥宴罢闻歌吹。\par
画毂走，香尘起。\par
冠压花枝驰万骑。\par
马行灯闹，凤楼帘卷，陆海龙鳌山对。\par
当年曾看天颜醉。\par
御杯举，欢声沸。\par
时节虽同悲乐异。\par
海风吹梦，岭猿啼月，一枕思归泪。\par

\section{西江月}
\textbf{惠洪}\par\vspace{0.5em}
入骨风流国色，透尘种性真香。\par
为谁风鬓涴新妆。\par
半树入村春暗。\par
雪压枝低篱落，月高影动池塘。\par
高情数笔寄微茫。\par
小寝初开雾帐。\par

\section{浪淘沙}
\textbf{惠洪}\par\vspace{0.5em}
城里久偷闲。\par
尘浣云衫。\par
此身已是再眠蚕。\par
隔岸有山归去好，万壑千岩。\par
霜晓更凭阑。\par
减尽晴岚。\par
微云生处是茅庵。\par
试问此生谁作伴，弥勒同龛。\par

\section{浪淘沙}
\textbf{惠洪}\par\vspace{0.5em}
山径晚樵远。\par
深壑孱颜。\par
孙山背后泊船看。\par
手把遗编披白帔，剩却清闲。\par
篱落竹丛寒。\par
渔业凋残。\par
水痕无底照秋宽。\par
好在夕阳凝睇处，数笔秋山。\par

\section{水调歌头}
\textbf{吕颐浩}\par\vspace{0.5em}
一片苍崖璞，孕秀自天锺。\par
浑如暖烟堆里，乍放力犹慵。\par
疑是犀眠海畔，贪玩烂银光彩，精魄入蟾宫。\par
泼墨阴云妒，蟾影淡朦胧。\par
沩山颂，戴生笔，写难穷。\par
些儿造化，凭谁细与问元工。\par
那用牧童鞭索，不入千群万队，扣角起雷同。\par
莫怪作诗手，偷入锦囊中。\par

\section{点绛唇}
\textbf{苏过}\par\vspace{0.5em}
新月娟娟，夜寒江静山衔斗。\par
起来搔首。\par
梅影横窗瘦。\par
好个霜天，闲却传杯手。\par
君知否。\par
乱鸦啼后。\par
归兴浓如酒。\par

\section{感皇恩}
\textbf{陆蕴}\par\vspace{0.5em}
残角两三声，催登古道。\par
远水长山又重到。\par
水声山色，看尽轮蹄昏晓。\par
风头日脚下，人空老。\par
匹马旧时，西征谈笑。\par
绿鬓朱颜正年少。\par
旗亭斗酒，任是十千倾倒。\par
面今酒兴减，诗情少。\par

\section{忆君王・忆王孙}
\textbf{谢克家}\par\vspace{0.5em}
依依宫柳拂宫墙。\par
楼殿无人春昼长。\par
燕子归来依旧忙。\par
忆君王。\par
月破黄昏人断肠。\par

\section{江神子・江城子}
\textbf{葛胜仲}\par\vspace{0.5em}
昏昏雪意惨云容。\par
猎霜风。\par
岁将穷。\par
流落天涯，憔悴一衰翁。\par
清夜小窗围兽火，倾酒绿，借颜红。\par
官梅疏艳小壶中。\par
暗香浓。\par
玉玲珑。\par
对景忽惊，身在大江东。\par
上国故人谁念我，晴嶂远，暮云重。\par

\section{蝶恋花}
\textbf{葛胜仲}\par\vspace{0.5em}
雨后春光浓似醉。\par
著柳催花，节物侵龙忌。\par
绣褓香阁当日佩。\par
紫兰宫堕人间世。\par
歌管停云香吐穗。\par
碧酒红裳，共祝鱼轩贵。\par
天上阿环金箓秘。\par
龟龄鹤寿三千岁。\par

\section{蝶恋花}
\textbf{葛胜仲}\par\vspace{0.5em}
共乐堂深帘不卷。\par
恻恻寒轻，二月春犹浅。\par
续寿竞来歌舞院。\par
龙涎香衬鲛绡段。\par
画栋朝飞双语燕。\par
端似知人，著意窥金盏。\par
柳外花前同祝愿。\par
朱颜长在年龄远。\par

\section{临江仙}
\textbf{葛胜仲}\par\vspace{0.5em}
郊外黄垓端可厌，归来移病香闺。\par
象床珍簟共委蛇。\par
耆婆寻草尽，天女散花迟。\par
小雨作寒秋意晚。\par
檐声与梦相宜。\par
冷侵罗幌酒烟微。\par
试评书五朵，何似画双眉。\par

\section{渔家傲}
\textbf{葛胜仲}\par\vspace{0.5em}
岩壑萦回云水窟。\par
林深路断迷烟客。\par
茅屋数椽携杖舄。\par
人寂寂。\par
侵檐万个琅玕碧。\par
倦客羁怀清似涤。\par
更无一点飞埃迹。\par
溪涨慢流过几席。\par
寒湜湜。\par
凫鷖点破琉璃色。\par

\section{渔家傲}
\textbf{葛胜仲}\par\vspace{0.5em}
叠叠云山供四顾。\par
簿书忙里偷闲去。\par
心远地偏陶令趣。\par
登览处。\par
清幽疑是斜川路。\par
野蔌溪毛供饮具。\par
此身甘被烟霞痼。\par
兴尽碧云催日暮。\par
招晚渡。\par
遥遥一叶随鸥鹭。\par

\section{鹧鸪天}
\textbf{葛胜仲}\par\vspace{0.5em}
小榭幽园翠箔垂。\par
云轻日薄淡秋晖。\par
菊英露浥渊明径，藕叶风吹叔宝池。\par
酬素景，泥芳卮。\par
老人痴钝强伸眉。\par
欢华莫遣笙歌散，归路从教灯影稀。\par

\section{鹧鸪天}
\textbf{葛胜仲}\par\vspace{0.5em}
婆律香深气味佳。\par
玻璃仙碗进流霞。\par
凝膏清涤高阳醉，灵液甘和正焙芽。\par
香染指，浪浮花。\par
加笾礼尽客还家。\par
贯珠声断红裳散，踏影人归素月斜。\par

\section{点绛唇}
\textbf{葛胜仲}\par\vspace{0.5em}
秋晚寒斋，藜床香篆横轻雾。\par
闲愁几许。\par
梦逐芭蕉雨。\par
云外哀鸿，似替幽人语。\par
归不去。\par
乱山无数。\par
斜日荒城鼓。\par

\section{行香子}
\textbf{葛胜仲}\par\vspace{0.5em}
风物飕飕。\par
木落沧洲。\par
渐老人、不奈悲秋。\par
羁怀都在，鬓上眉头。\par
似休文瘦，文通恨，子山愁。\par
庭梧影薄，篱菊香浮。\par
强招寻、聊命朋俦。\par
穷通皆梦，今古如流。\par
且渊明径，子猷舫，仲宣楼。\par

\section{诉衷情}
\textbf{葛胜仲}\par\vspace{0.5em}
清明寒食景暄妍。\par
花映碧罗天。\par
参差捍拨齐奏，丰颊拥芳筵。\par
逢诞日，揖真仙。\par
托炉烟。\par
朱颜长似，头上花枝，岁岁年年。\par

\section{水调歌头}
\textbf{葛胜仲}\par\vspace{0.5em}
夜泛南溪月，光影冷涵空。\par
棹飞穿碎金电，翻动水精宫。\par
横管何妨三弄，重醑仍须一斗，知费几青铜。\par
坐久桂花落，襟袖觉香浓。\par
庾公阁，子猷舫，兴应同。\par
从来好景良夜，我辈敢情锺。\par
但恐仙娥川后，嫌我尘容俗状，清境不相容。\par
击汰同情赏，赖有紫溪翁。\par

\section{水调歌头}
\textbf{葛胜仲}\par\vspace{0.5em}
下濑惊船驶，挥麈恐尊空。\par
谁吹尺八寥亮，嚼徵更含宫。\par
坐爱金波潋滟，影落蒲萄涨绿，夜漏尽移铜。\par
回棹携红袖，一水带香浓。\par
坐中客，驰隽辩，语无同。\par
青鞋黄帽，此乐谁肯换千锺。\par
岩壑从来无主，风月故应长在，赏不待先容。\par
羽化寻烟客，家有左仙翁。\par

\section{水调歌头}
\textbf{葛胜仲}\par\vspace{0.5em}
胜友欣倾盖，羁宦懒书空。\par
爱君笔力清壮，名已在蟾宫。\par
萧散英姿直上，自有练裙葛帔，岂待半通铜。\par
长短作新语，墨纸似鸦浓。\par
山吐月，溪泛艇，率君同。\par
吾侪轰饮文字，乐不在歌锺。\par
今夜长风万里，且倩泓澄浩荡，一为洗尘容。\par
世上闲荣辱，都付寒边翁。\par

\section{木兰花・玉楼春}
\textbf{葛胜仲}\par\vspace{0.5em}
木阑干外池光阔。\par
午夜乔林迷岸樾。\par
掠船凉吹起青苹，萦水歌声欺白雪。\par
檀郎响趁红牙节。\par
胡语嘈嘈仍切切。\par
人生何乐似同襟，莫待骊驹声惨咽。\par

\section{满庭霜・满庭芳}
\textbf{葛胜仲}\par\vspace{0.5em}
百不为多，一不为少，阿谁昔仕吾邦。\par
共推任笔，洪鼎力能扛。\par
不为桃花禄米，雠书倦、一苇横江。\par
招寻处，徒行曳杖，曾不拥麾幢。\par
山川，真大好，鱼矶无恙，密岭难双。\par
听讼诉多就，樵坞僧窗。\par
岁月音容远矣，风流在、遐想心降。\par
云烟路，搜奇吊古，时为酹空缸。\par

\section{醉蓬莱}
\textbf{葛胜仲}\par\vspace{0.5em}
望葱葱佳气，虹渚祥开，斗枢光绕。\par
析木天津，正灵晖腾照。\par
鹭缀分班，象胥交贡，奉御觞清晓。\par
玉殿寒轻，金徒漏永，瑞炉烟袅。\par
万宇均欢，示慈颁燕，寿祝南山，庆均凫藻。\par
缥缈红云，望九重天表。\par
舞兽锵洋，抃鳌欣戴，度管弦声杳。\par
历草长新，蟠桃永秀，与天难老。\par

\section{西江月}
\textbf{葛胜仲}\par\vspace{0.5em}
羁宦新来作恶，穷途谁肯相从。\par
追攀十日水云中。\par
情谊知君独重。\par
寂寂回廊小院，冥冥细雨尖风。\par
凤山香雪定应空。\par
昨夜疏枝入梦。\par

\section{西江月}
\textbf{葛胜仲}\par\vspace{0.5em}
山镇红桃阡陌，烟迷绿水人家。\par
尘容误到只惊嗟。\par
骨冷玉堂今夜。\par
莫对佳人锦瑟，休辞洞府流霞。\par
峰回路转乱云遮。\par
归去空传图画。\par

\section{南乡子}
\textbf{葛胜仲}\par\vspace{0.5em}
柳岸正飞绵。\par
选胜斋轻漾碧涟。\par
笑语忘怀机事尽，鸥边。\par
万顷溪光上下天。\par
菰苇久延缘。\par
不觉遥峰霭暮烟。\par
对酒莫嫌红粉陋，婵娟。\par
自有孤高月妇仙。\par

\section{浣溪沙}
\textbf{葛胜仲}\par\vspace{0.5em}
可惜随风面旋飘。\par
直须烧烛看娇娆。\par
人间花月更无妖。\par
浓丽独将春色殿，繁华端合众芳朝。\par
南床应为醉陶陶。\par

\section{浣溪沙}
\textbf{葛胜仲}\par\vspace{0.5em}
通白轻红溢万枝。\par
浓香百和透丰肌。\par
丹山威凤势将飞。\par
玉镜台前呈国艳，沈香亭北映朝曦。\par
如花惟有上皇妃。\par

\section{浣溪沙}
\textbf{葛胜仲}\par\vspace{0.5em}
斗鸭栏边晓路沾。\par
华堂醉赏轴珠帘。\par
插花人好手纤纤。\par
遮护轻寒施翠幄，标题仙品露牙签。\par
词人遗恨独江淹。\par

\section{西江月}
\textbf{葛胜仲}\par\vspace{0.5em}
山外半规残日，云边一缕馀霞。\par
满城飞雪散苕花。\par
万顷溪连罨画。\par
柳恽风流旧国，鹤龄潇洒人家。\par
肯嗟流落在天涯。\par
云水从今起价。\par

\section{西江月}
\textbf{葛胜仲}\par\vspace{0.5em}
晚路交游绿酒，平生志趣青霞。\par
霜风时节近黄花。\par
泛宅舟将鹢画。\par
不分两溪明月，夜深只属渔家。\par
今朝清赏寄情涯。\par
肯向萦涂索价。\par

\section{蝶恋花}
\textbf{葛胜仲}\par\vspace{0.5em}
风过涟漪纹縠细。\par
十指香檀，惊破交禽睡。\par
野蔌溪毛真易致。\par
风流未减兰亭会。\par
击汰千艘供洛禊。\par
映水垂杨，万缕拖浓翠。\par
小海一声波上戏。\par
殷勤留客千金意。\par

\section{临江仙}
\textbf{葛胜仲}\par\vspace{0.5em}
自古吴兴称冷僻，菰城水浸粼粼。\par
回星难望使车尘。\par
如何三日饮，并有五行人。\par
文似枚皋加敏速，记书易若张巡。\par
幕中无用郄嘉宾。\par
他年浮枣会，莫忘两溪春。\par

\section{临江仙}
\textbf{葛胜仲}\par\vspace{0.5em}
千古乌程新酿美，玉觞风过粼粼。\par
歌声未办起梁尘。\par
九天持斧客，来作绣衣人。\par
夙有辞华惊乙览，传闻献颂东巡。\par
未应握节久宾宾。\par
一封驰诏旨，却醉上林春。\par

\section{临江仙}
\textbf{葛胜仲}\par\vspace{0.5em}
小样洪河分九曲，飞泉环绕粼粼。\par
青莲往事已成尘。\par
羽觞浮玉甃，宝剑捧金人。\par
绿绮且依流水调，蓬蓬釂鼓催巡。\par
玉堂词客是佳宾。\par
茂林修竹地，大胜永和春。\par

\section{定风波}
\textbf{葛胜仲}\par\vspace{0.5em}
千叠云山万里流。\par
坐中碧落与鳌头。\par
真意见嬉吾已领。\par
烟景。\par
不辞捧诏久汀洲。\par
老去一官真是漫。\par
溪岸。\par
独馀此兴未能收。\par
留与吴儿传胜事。\par
长记。\par
赤栏桥上揽清秋。\par

\section{定风波}
\textbf{葛胜仲}\par\vspace{0.5em}
共喜新凉大火流。\par
一声水调听歌头。\par
况有修蛾兼粉领。\par
佳景。\par
谢公无不碍沦洲。\par
平昔短檠真大漫。\par
气岸。\par
老来都向酒杯收。\par
云水光中修禊事。\par
犹记。\par
转头不觉已三秋。\par

\section{浣溪沙}
\textbf{葛胜仲}\par\vspace{0.5em}
今夜风光恋渚苹。\par
欲教四角出车轮。\par
金钗离立座生春。\par
神女恍惊巫峡梦，飞琼原是阆风人。\par
诏封后院宠儒臣。\par

\section{浣溪沙}
\textbf{葛胜仲}\par\vspace{0.5em}
溪岸沈深属泛苹。\par
倾城容貌此推轮。\par
可怜虚度二年春。\par
暮暮来时骚客赋，朝朝新处后庭人。\par
天留花月伴羁臣。\par

\section{浣溪沙}
\textbf{葛胜仲}\par\vspace{0.5em}
东道殷勤玉斝飞。\par
华灯倾国拥珠玑。\par
玉奴嫌瘦玉环肥。\par
缥缈幸闻缑岭曲。\par
参差犹隔夏侯衣。\par
放开云月出清辉。\par

\section{瑞鹧鸪}
\textbf{葛胜仲}\par\vspace{0.5em}
两年人住岂无情。\par
别乘辞华四水清。\par
何事千锺勤饮饯，故知一别未能轻。\par
解龟虽幸樊笼出，挂席还愁海汐平。\par
江草江花都是泪，骊驹休作断肠声。\par

\section{浪淘沙}
\textbf{葛胜仲}\par\vspace{0.5em}
步屧对东风。\par
细探春工。\par
百花堂下牡丹丛。\par
莫恨使君来便去，不见鞓红。\par
雾眼一衰翁。\par
无意芳秾。\par
年来结习已成空。\par
寄语国香雕槛里，好为人容。\par

\section{蓦山溪}
\textbf{葛胜仲}\par\vspace{0.5em}
望云门外。\par
油壁如流水。\par
空巷逐朱幡，步春风、香河七里。\par
冶容炫服，摸石道宜男，穿翠霭，度飞桥，影在清漪里。\par
秦头楚尾。\par
千古风流地。\par
试问汉江边，有解佩、行云旧事。\par
主人是客，一笑强颁春，烧灯后，赏花前，遥忆年年醉。\par

\section{蓦山溪}
\textbf{葛胜仲}\par\vspace{0.5em}
春风野外。\par
卵色天如水。\par
鱼戏舞绡纹，似出听、新声北里。\par
追风骏足，千骑卷高冈，一箭过，万人呼，雁落寒空里。\par
天穿过了，此日名穿地。\par
横石俯清波。\par
竞追随、新年乐事。\par
谁怜老子，使得暂遨游，争捧手，乍凭肩，夹道游人醉。\par

\section{减字木兰花}
\textbf{吴则礼}\par\vspace{0.5em}
团团璧月。\par
今夜广寒真秀发。\par
何处吹笙。\par
催得清霜满凤城。\par
淮南好梦。\par
镜里星星还种种。\par
犹记银床。\par
曾为凉洲唤玉觞。\par

\section{减字木兰花}
\textbf{吴则礼}\par\vspace{0.5em}
淮天不断。\par
点缀南云秋几雁。\par
白露沾衣。\par
始是银屏梦觉时。\par
别离怀抱。\par
消得镜中青鬓老。\par
小字无能。\par
烦寄平安一纸书。\par

\section{减字木兰花}
\textbf{吴则礼}\par\vspace{0.5em}
斑斑小雨。\par
初入高梧黄叶暮。\par
又是重阳。\par
昨夜西风作许凉。\par
鲜鲜丛菊。\par
只解凋人双鬓绿。\par
试傍清尊。\par
分付幽香与断魂。\par

\section{减字木兰花}
\textbf{吴则礼}\par\vspace{0.5em}
九年离别。\par
梦里相逢端怕说。\par
携手河梁。\par
雁噭淮天如许长。\par
鲈鱼正美。\par
白发季鹰聊启齿。\par
后夜江干。\par
与把梅花子细看。\par

\section{减字木兰花}
\textbf{吴则礼}\par\vspace{0.5em}
淮山清夜。\par
镜面平铺纤月挂。\par
端是生还。\par
同倚西风十二栏。\par
休论往事。\par
投老相逢真梦寐。\par
两鬓疏疏。\par
好在松江一尺鲈。\par

\section{满庭芳}
\textbf{吴则礼}\par\vspace{0.5em}
声促铜壶，灰飞玉琯，梦惊偷换年华。\par
江南芳信，疏影月横斜。\par
又喜椒觞到手，宝胜里、仍翦金花。\par
钗头燕，妆台弄粉，梅额故相夸。\par
隼旟，人未老，东风袅袅，已傍高牙。\par
渐园林月永，叠鼓凝笳。\par
小字新传秀句，歌扇底、深把流霞。\par
聊行乐，他时画省，归近紫皇家。\par

\section{满庭芳}
\textbf{吴则礼}\par\vspace{0.5em}
玉垒尊罍，清秋关塞，正宜催唤香醪。\par
凉风吹帽，横槊试登高。\par
相见征西旧事，龙山会、宾主俱豪。\par
来群雁，凉生画角，红叶聚亭皋。\par
南洲、应好在，长关梦眼，天际云涛。\par
有一簪黄菊，两鬓霜毛。\par
快把金荷共倒，凝望久、只遣魂消。\par
君须听，新翻燕乐，馀韵响檀槽。\par

\section{木兰花慢}
\textbf{吴则礼}\par\vspace{0.5em}
尽晴春自老，乍翠巘、出清流。\par
望杳杳飞旌，翩翩戍骑，初过边头。\par
幽花尚猗短岸，渐鸣禽、唤友绕行輈。\par
端有雕戈锦领，竟驰騕袅骅骝。\par
凝眸。\par
雁入长天，羌管罢、陇云愁。\par
共解鞍临水，雷惊电散，雪溅霜浮。\par
玉觞正风味好，对幽香、堕蕊且消忧。\par
莫以纶巾羽扇，便忘绿浦沧洲。\par

\section{醉落魄・一斛珠}
\textbf{吴则礼}\par\vspace{0.5em}
梅花褪雪。\par
赏心莫使慵欢悦。\par
大家且恁同攀折。\par
馀蕊残英，偏称淡笼月。\par
当初相见花初发。\par
如今花谢人离缺。\par
一年又比一年别。\par
惟有花枝，只似旧时节。\par

\section{踏莎行}
\textbf{吴则礼}\par\vspace{0.5em}
一片花飞，青春已减。\par
可堪南陌红千点。\par
生憎杨柳要藏鸦，东风只遣横笳怨。\par
看定新巢，初怜语燕。\par
游丝正把残英罥。\par
酒尊也会不相违，风光本自同流转。\par

\section{鹧鸪天}
\textbf{吴则礼}\par\vspace{0.5em}
永遇英雄际会时。\par
垂天鹏翼逐云飞。\par
退朝日上青花道，催直霜零赤雁池。\par
鸣汉履，侍唐眉。\par
渭川莘野晚追随。\par
归来仍对金銮老，三峡词源气未衰。\par

\section{鹧鸪天}
\textbf{吴则礼}\par\vspace{0.5em}
作赋丁年厌兔园。\par
紫微深锁九重关。\par
花墩屡赐清闲燕，文石难忘咫尺颜。\par
烟外屐，水边山。\par
纶巾羽扇五湖间。\par
自怜季子貂裘敝，来与机云相对闲。\par

\section{鹧鸪天}
\textbf{吴则礼}\par\vspace{0.5em}
衮绣三朝社稷臣。\par
旧调元鼎斡洪钧。\par
垂绅屡转龙墀日，接膝潜回黼座春。\par
金凿落，玉麒麟。\par
凤鸣良月庆佳辰。\par
巨鳌行听扃华禁，又起商周梦卜人。\par

\section{红楼慢}
\textbf{吴则礼}\par\vspace{0.5em}
声慑燕然，势压横山，镇西名重榆塞。\par
干霄百雉朱阑下，极目长河如带。\par
玉垒凉生过雨，帘卷晴岚凝黛。\par
有城头、钟鼓连云，殷春雷天外。\par
长啸，畴昔驰边骑。\par
听陇底鸣笳，风搴双旆。\par
霜髯飞将曾百战，欲掳名王朝帝。\par
锦带吴钩未解，谁识凭栏深意。\par
空沙场，牧马萧萧晚无际。\par

\section{声声慢}
\textbf{吴则礼}\par\vspace{0.5em}
林塘朱夏，雨过斑斑，绿苔绕地初遍。\par
叶底雏莺，犹记日斜春晚。\par
芙蕖靓妆红粉，傍高荷、闲倚歌扇。\par
轻风起，縠纹滟滟，翠生波面。\par
可是追凉月下，清坐久，微云屡遮星汉。\par
露湿纶巾，遥望玉清台殿。\par
白头共论胜事，须偿五湖深愿。\par
南枝好，有南飞乌鹊，绕枝低转。\par

\section{水龙吟}
\textbf{吴则礼}\par\vspace{0.5em}
秋生泽国，无边落木，又作萧萧下。\par
澄江过雨，凉飙吹面，黄花初把。\par
苍鬓羁孤，粗营鸡黍，浊醪催贳。\par
对斜斜露脚，寒香正好，幽人去、空惊咤。\par
头上纶巾醉堕，要猗眠、水云萦舍。\par
牵衣儿女，归来欢笑，仍邀同社。\par
月底蓬门，一株江树，悲虫鸣夜。\par
把茱萸细看，牛山底事，强成沾洒。\par

\section{眼儿媚}
\textbf{李德载}\par\vspace{0.5em}
雪儿魂在水云乡。\par
犹忆学梅妆。\par
玻璃枝上，体薰山麝，色带飞霜。\par
水边竹外愁多少，不断俗人肠。\par
如何伴我，黄昏携手，步月斜廊。\par

\section{早梅芳近・早梅芳}
\textbf{李德载}\par\vspace{0.5em}
深院静，小阑傍。\par
标致不寻常。\par
尽他桃杏占风光。\par
谁敢斗新妆。\par
玉堂中，梁苑里。\par
休把雪来轻比。\par
莫吹长笛巧催残。\par
留取月中看。\par

\section{早梅芳近・早梅芳}
\textbf{李德载}\par\vspace{0.5em}
残腊里，早梅芳。\par
春信报新阳。\par
晓来枝上斗寒光。\par
轻点寿阳妆。\par
雪难欺，霜莫妒。\par
别是一般风措。\par
望林人意正夭饶。\par
又看长新条。\par

\section{鹧鸪天}
\textbf{赵子发}\par\vspace{0.5em}
约略应飞白玉盘。\par
明楼渐放满轮寒。\par
天垂万丈清光外，人在三秋爽气间。\par
闻叶吹，想风鬟。\par
浮空仿佛女乘鸾。\par
此时不合人间有，尽入嵩山静夜看。\par

\section{洞仙歌}
\textbf{赵子发}\par\vspace{0.5em}
荒山明月，下有云来去。\par
深夜纤毫静可数。\par
问古今底事，留此空光，修月户、犹是当年玉斧。\par
思君持羽扇，来伴微吟，水佩风环饮松露。\par
待勾漏丹成，约与轻飞，人间世、不知归处。\par
更长啸、馀声振林溪，见乱红惊飞，半岩花雨。\par

\end{document}